\documentclass[12pt]{article}
\usepackage[spanish]{babel}
\usepackage[T1]{fontenc}
\usepackage[utf8]{inputenc}
\usepackage{lmodern}
\usepackage[margin=1in]{geometry}
\usepackage{xcolor}
\usepackage{soulutf8}
\usepackage{enumitem}
\usepackage{setspace}

\setlength{\parindent}{0pt}
\setlength{\parskip}{0.4em}
\setstretch{1.08}

\definecolor{hlblue}{RGB}{219,244,255}
\definecolor{hlpink}{RGB}{248,214,228}
\definecolor{hlgreen}{RGB}{217,247,208}
\definecolor{hlyellow}{RGB}{247,240,170}
\definecolor{hlorange}{RGB}{249,220,179}

\newcommand{\hlbluebox}[1]{{\sethlcolor{hlblue}\hl{#1}}}
\newcommand{\hlpinkbox}[1]{{\sethlcolor{hlpink}\hl{#1}}}
\newcommand{\hlgreenbox}[1]{{\sethlcolor{hlgreen}\hl{#1}}}
\newcommand{\hlyellowbox}[1]{{\sethlcolor{hlyellow}\hl{#1}}}
\newcommand{\hlorangebox}[1]{{\sethlcolor{hlorange}\hl{#1}}}

\begin{document}

{\Large\textbf{Ejemplo de Sistematización}}

\vspace{1em}

{\large\textbf{Diagnóstico}}

\vspace{0.6em}

\textbf{LÍNEA DE INVESTIGACIÓN:} Inteligencia Artificial y Comunicación Inalámbrica.

En los \hlpinkbox{museos y espacios culturales}, el \hlgreenbox{servicio de guía virtual} apoyado en \hlbluebox{aplicaciones móviles e IoT con señales Wi-Fi/BLE} presenta limitaciones para ofrecer información contextualizada y oportuna al visitante, debido a \hlpinkbox{la baja precisión de localización indoor}, \hlpinkbox{la heterogeneidad de dispositivos}, \hlpinkbox{la inestabilidad del RSSI}, \hlpinkbox{el efecto multipath} y \hlpinkbox{las variaciones temporales del entorno inalámbrico}. Esta situación dificulta identificar de manera confiable la sala, zona u objeto próximo al usuario, afecta la activación correcta de la guía hablada y reduce \hlyellowbox{la calidad del servicio de guía virtual} dentro del recorrido museográfico.

\vspace{0.8em}

\textbf{Tema General:}

\hlpinkbox{Museos y espacios culturales}

\vspace{0.8em}

\textbf{Temas Particulares:}

\begin{enumerate}[leftmargin=2.2em]
  \item \hlgreenbox{Servicios de guía virtual} en \hlpinkbox{museos y espacios culturales}
  \item \hlgreenbox{Servicios de guía virtual} con \hlbluebox{aplicaciones móviles e IoT con señales Wi-Fi/BLE} en \hlpinkbox{museos y espacios culturales}
  \item \hlgreenbox{Servicios de guía virtual} con detección de proximidad y \hlpinkbox{localización indoor} en \hlpinkbox{museos y espacios culturales}
\end{enumerate}

\vspace{0.6em}

\textbf{Temas Específicos:}

\begin{enumerate}[leftmargin=2.2em]
  \item \hlyellowbox{Baja calidad del servicio de guía virtual} en \hlgreenbox{servicios de guía virtual} con \hlbluebox{aplicaciones móviles e IoT con señales Wi-Fi/BLE} en \hlpinkbox{museos y espacios culturales}
  \item \hlpinkbox{Baja precisión de localización indoor} en \hlgreenbox{servicios de guía virtual} con \hlbluebox{aplicaciones móviles e IoT con señales Wi-Fi/BLE} en \hlpinkbox{museos y espacios culturales}
  \item \hlpinkbox{Alta sensibilidad a la heterogeneidad de dispositivos, inestabilidad del RSSI, efecto multipath y variaciones temporales} en \hlgreenbox{servicios de guía virtual} con \hlbluebox{aplicaciones móviles e IoT con señales Wi-Fi/BLE}
\end{enumerate}

\textbf{Problema de Estudio:}

\hlyellowbox{Baja calidad del servicio de guía virtual} en \hlgreenbox{servicios de guía virtual} con \hlbluebox{aplicaciones móviles e IoT con señales Wi-Fi/BLE} en \hlpinkbox{museos y espacios culturales}, \hlbluebox{2026}

\vspace{0.6em}

\textbf{Título Preliminar de la Tesis:}

\newpage

\hlyellowbox{Calidad del servicio de guía virtual} en \hlgreenbox{servicios de guía virtual} con \hlbluebox{aplicaciones móviles e IoT con señales Wi-Fi/BLE} en \hlpinkbox{museos y espacios culturales}, \hlbluebox{2026}

\vspace{0.9em}

\textbf{Antecedentes Bibliográficos}

Gufran, D., Tiku, S., \& Pasricha, S. (2023). \textit{STELLAR: Siamese Multiheaded Attention Neural Networks for Overcoming Temporal Variations and Device Heterogeneity With} \hlgreenbox{\textit{Indoor Localization}}. \textit{IEEE Journal of Indoor and Seamless Positioning and Navigation}.

Singampalli, A., Sri Lakshmi, G. R. R., \& Pasricha, S. (2024). \textit{xKAN: Explainable AI for Smartphone Based} \hlgreenbox{\textit{Indoor Localization}} \textit{using} \hlbluebox{\textit{Wi-Fi Fingerprinting}}. \textit{IEEE Access}.

Tiku, S., Gufran, D., \& Pasricha, S. (2022). \textit{Multi-Head Attention Neural Network for Smartphone Invariant} \hlgreenbox{\textit{Indoor Localization}}. En \textit{International Conference on Indoor Positioning and Indoor Navigation (IPIN)}.

Supanakoon, P., \& Promwong, S. (2021). \textit{Evaluation of Distance Error with Bluetooth} \hlbluebox{\textit{Low Energy Transmission Model}} \textit{for} \hlgreenbox{\textit{Indoor Positioning}}. \textit{Journal of Mobile Multimedia}, 17(4), 707--722.

Xu, S., Chen, R., Guo, G., Li, Z., Qian, L., Ye, F., Liu, Z., \& Huang, L. (2021). \hlbluebox{\textit{Bluetooth, Floor-Plan, and Microelectromechanical Systems-Assisted Wide-Area Audio}} \hlgreenbox{\textit{Indoor Localization}} \textit{System: Apply to Smartphones}. \textit{IEEE Transactions on Industrial Electronics}.

Tsuchiya, Y., Suga, N., Uruma, K., \& Fujisawa, M. (2024). \textit{Integration of CSI and Camera Image for} \hlgreenbox{\textit{Indoor Wireless Localization}}. \textit{IEEE Access}.

Yaman, I., Tian, G., Tegler, E., Gulin, J., Challa, N., Tufvesson, F., Edfors, O., Åström, K., Malkowsky, S., \& Liu, L. (2024). \textit{LuViRA Dataset Validation and Discussion: Comparing Vision, Radio, and Audio Sensors for} \hlgreenbox{\textit{Indoor Localization}}. \textit{IEEE Journal of Indoor and Seamless Positioning and Navigation}.

Salamah, A. H., Tamazin, M., Sharkas, M. A., \& Khedr, M. (2016). \textit{An Enhanced} \hlbluebox{\textit{WiFi Indoor Localization System}} \textit{Based on Machine Learning}. En \textit{International Conference on Indoor Positioning and Indoor Navigation (IPIN)}.

Zhang, J., Sun, J., Wang, H., Xiao, W., \& Tan, L. (2017). \textit{Large-Scale} \hlbluebox{\textit{WiFi Indoor Localization}} \textit{via Extreme Learning Machine}. En \textit{Proceedings of the 36th Chinese Control Conference}.

Aguilar Herrera, J. C., Plöger, P. G., Hinkenjann, A., Maiero, J., Flores, M., \& Ramos, A. (2014). \textit{Pedestrian} \hlgreenbox{\textit{Indoor Positioning}} \textit{Using Smartphone Multi-sensing, Radio Beacons, User Positions Probability Map and IndoorOSM Floor Plan Representation}. En \textit{International Conference on Indoor Positioning and Indoor Navigation (IPIN)}.

\vspace{0.8em}

\textbf{Título Tentativo de la Tesis:}

\newpage

\hlorangebox{Modelo de aprendizaje profundo con mecanismos de atención} para mejorar \hlyellowbox{la calidad del servicio de guía virtual} en \hlgreenbox{servicios de guía virtual} con \hlbluebox{aplicaciones móviles e IoT con señales Wi-Fi/BLE} en \hlpinkbox{museos y espacios culturales}, \hlbluebox{2026}

\vspace{1em}

\textbf{Variable Independiente:} \hlorangebox{Modelo de aprendizaje profundo con mecanismos de atención}

\textbf{Variable Dependiente:} \hlyellowbox{Calidad del servicio de guía virtual}

\textbf{Objeto de Estudio:} \hlgreenbox{Servicios de guía virtual}

\textbf{Alcance Espacial:} \hlpinkbox{Museos y espacios culturales}

\textbf{Alcance Temporal:} \hlbluebox{2026}

\end{document}
